\documentclass[Comic Sans MS]{themain}
\begin{document}
% Ordre des compétences sur Pronote
% Transversale
% Chercher
% Modeliser
% Représenter
% Raisonner
% Calculer
% Communiquer
\level{6e}
\sequence{Aborder\ la\ g\'eometrie}\hiddentrue
    \addEvalCmp{Je connais les objets géométriques usuels.&représenter}
    \addEvalCmp{Je comprends les notations géométriques.&communiquer}
    \addEvalCmp{Je présente correctement ma(mes) feuille(s), mon écriture est lisible.&transversale}
\cours
\activites{
    notation_geometrie-1,
    appartenance-1,
    programme_de_trace-1,
    programme_de_trace-2%
    }
\begin{exercices}
    \doublesheet{0361-c,0362-c}{0363-c,0364-c}
    \simplesheet{0365-c,0366-c}
    \doublesheet{0368-c,0369-c}{0367-c,0370-c}
    \simplesheet{0142-c,0143-c}
\end{exercices}
\questionsFlash{
    tables_de_multiplication-2,
    tables_de_multiplication-3,
    tables_de_multiplication-1,
    notations_geometrie-1,
    completer_liste_logique-1,
    completer_liste_logique_difficile-1%
}
\setEvalCmp{\cmpItem0\cmpItem1\cmpItem2}
\evaluation{% 
    format=DM,
    duree=À rendre le 14/09/2022,
    consignes={Répondre directement sur la feuille.}
}{0135-c/5,0136-a/5}
\evaluationAB{%
    duree=50 min,
    consignes={Répondre directement sur la feuille.},
    skip page id=0135%
}{0140/6,0135/5,0141/3,0142/3,0143/3}
\evaluation{%
    format=DS,
    sujet=C,
    duree=50 min,
    consignes={Répondre directement sur la feuille.},
    font=sf,
    skip page id=0135-b%
}{0140-b/6,0135-b/5,0141-b/6,0143-b/3}
\endsequence
\sequence{Entiers\ naturels}\hiddentrue
	\addEvalCmp{Je sais poser et effectuer ces trois opérations&calculer}
	\addEvalCmp{Je conclue à l'aide d'une phrase réponse&communiquer}
	\addEvalCmp{Je présente correctement ma(mes) feuille(s), mon écriture est lisible&transversale}
\cours
\activites{
    ecriture_nombres_entiers-1%
}
\begin{exercices}
    \doublesheet{0151-c}{0153-c,0154-c}
    \doublesheet{0403-c,0404-c}{0405-c,0406-c}
\end{exercices}
\evaluationAB{%
    duree=50 min,
    consignes={Pour chaque problème, répondre par au moins un
    calcul et une phrase réponse.},
    calculatrice=interdite,
    skip page id=0156%
}{0587/4,0147/2,0148/2,0149/2,0151/4,0156/3,0403/3,0152/ 1 Point bonus}
\evaluation{%
    format=DS,
    sujet=C,
    duree=50 min,
    font=sf,
    consignes={Pour chaque problème, répondre par au moins un
    calcul et une phrase réponse.},
    calculatrice=autorisée%
}{0587-a/3,0147-a/4,0148-a/4,0149-a/4,0403-a/2,0151-a/3}
\endsequence
\sequence{Angles}\hiddentrue
    \addEvalCmp{Je présente correctement ma(mes) feuille(s), mon écriture est lisible.&transversale}
    \addEvalCmp{Je connais le vocabulaire associé aux angles.&transversale}
    \addEvalCmp{Je sais tracer ou mesurer un angle.&représenter}
\cours
\activites{
    angle-1,
    notation_angles-1,
    nature_angles-1,
    mesurer_angle_partie-1,
    mesurer_angle_partie-2,
    tracer_un_angle-1%
    }
\begin{exercices}
    \doublesheet{0165-c,0166-c}{0167-c}
    \doublesheet{0168-c,0169-c}{0170-c}
    \doublesheet{0173-c}{0172-c,0171-c}
    \doublesheet{0427-c,0427-d}{0428-c,0428-d,0429-c}
    \doublesheet{0176-c,0177-c}{0178-c}
    \simplesheet{0188-c}
\end{exercices}
\questionsFlash{
    calculs_astucieux_entiers-2,
    le_compte_est_bon_simple-1,
    notations_geometrie-2,
    le_compte_est_bon_simple-3,
    nature_des_angles-1%
}
\setEvalCmp{\cmpItem0\cmpItem1\cmpItem2}
\evaluationAB{
    duree=50 min,
    consignes={Répondre par des {\bfseries phrases}. Répondre sur
    le sujet uniquement pour la question {\bfseries 1)} de l'exercice 3.},
    skip page id=0173%
}{0167/4,0173/3,0190/5,0428/3,0188/2,0189/3}
\evaluation{
    format=DS,
    sujet=C,
    duree=50 min,
    font=sf,
    consignes={Répondre par des {\bfseries phrases}. Répondre sur
    le sujet uniquement pour la question {\bfseries 1)} de l'exercice 3.},
    skip page id=0173-a%
}{0167-a/6,0173-a/6,0190-a/5,0428-a/3}
\endsequence
\sequence{Fractions}\hiddentrue
	\addEvalCmp{Je présente correctement ma(mes) feuille(s), mon écriture est lisible&transversale}
	\addEvalCmp{Je reconnais les différentes représentations d'un même nombre&représenter}
	\addEvalCmp{Je suis capable d'appliquer un programme de tracé&représenter}
\cours
\activites{
    fractions_et_partages-1,
    fractions_et_partages-2,
    fractions_droite_graduee-1,
    fractions_decimales-1,
    ecriture_decimale-1,
    zeros_inutiles-1%
}
\begin{exercices}
    \doublesheet{0197-c,0198-c}{0199-c}
    \doublesheet{0193-c}{0193-d}
    \doublesheet{0194-c}{0195-c,0196-c}
    \doublesheet{0222-c}{0223-c}
\end{exercices}
\questionsFlash{
    multiplications_astucieuses_entiers-1,
    multiplications_astucieuses_entiers-2,
    multiplications_astucieuses_entiers-3,
    comparaison_fractions-1,
    nature_des_angles-1,
    vocabulaire_operations-2%
}
\evaluationAB{
    duree=50 min,
    consignes={\noindent Répondre sur votre feuille pour les
    exercices 1 et 2. Pour l'exercice 3, vous pouvez vous aider d'un
    tableau que vous tracerez sur le sujet.},
    calculatrice=interdite,
    first group=\multicolstrue
}{{0198/2,0222/3,0194/6,0485/1,0486/1,0487/2},{0488/5}}
\evaluation{
    sujet=C,
    duree=50 min,
    consignes={\noindent Répondre sur votre feuille pour les
    exercices 1 et 2. Pour l'exercice 3, vous pouvez vous aider d'un
    tableau que vous tracerez sur le sujet.},
    calculatrice=autorisée,
    font=sf%
}{0198-a/2,0222-a/3,0194-a/6,0485-a/1,0486-a/1,0487-a/2,0488-a/5}
\endsequence
\sequence{Droites\ perpendiculaires\ et\ parall\`eles}\hiddentrue
	\addEvalCmp{Je connais les notations géométriques&communiquer}
	\addEvalCmp{J'utilise mes outils de géométrie avec précision&représenter}
	\addEvalCmp{Je reconnais les différentes représentations d'un même nombre&représenter}
\cours
\activites{
    perpendicularite-1,
    propriete_paralleles-1%
}
\begin{exercices}
    \simplesheet{0588-c}
    \doublesheet{0589-c}{0590-c}
    \simplesheet{0593-c}
    \doublesheet{0591-c}{0592-c}
\end{exercices}
\questionsFlash{%
    fraction_decimale_en_ecriture_decimale-1%
}
\evaluation{
    format=DS,
    sujet=A,
    duree=50 min,
    calculatrice=interdite,
    font=sf,
    consignes={Toutes les constructions doivent être réalisées sur le
    sujet en utilisant les outils de géométrie.}%
}{0593-c/4,0591-c/5,0594-c/5,0194-d/6}
\endsequence
\sequence{Op\'erations\ sur\ les\ nombres\ d\'ecimaux}\hiddentrue
\cours
\activites{
    multiplication_decimaux-1,
    multiplication_mental_decimaux-1,
    multiplier_par_01_001-1,
    division_decimale-1%
    }
\begin{exercices}
    \doublesheet{0242-c,0242-d,0243-c}{0244-c,0245-c}
    \doublesheet{0007-c,0008-d,0009-a}{0010-a,0006-d}
    \doublesheet{0253-c,0254-c,0256-c}{0255-c}
    \doublesheet{0514-c}{0515-c,0516-c}
    \doublesheet{0517-c}{0518-c}
    \doublesheet{0519-c}{0520-c}
    \doublesheet{0274-c,0275-c,0276-c}{0277-c,0278-c,0279-c}
\end{exercices}
\questionsFlash{%
    le_compte_est_bon_simple-2,
    multiplier_par_05_et_025-1,
    multiplier_par_05_et_025-2,
    comparaison_fractions-2,
    multiplier_decimaux-1,
    multiplier_par_05_et_025-5,
    multiplication_division_par_10_100_1000,
    multiplication_division_par_10_100_1000-2%
}
\evaluationAB{
    duree=50 min,
    consignes={Ne pas répondre sur le sujet.},
    calculatrice = interdite%
}{0525/4,0526/2,0008/4,0253/4,0552/2,0277/2,0006/2}
\evaluation{
    format=DS,
    sujet=C,
    font=sf,
    duree = 50 min,
    consignes={Ne pas répondre sur le sujet.},
    calculatrice = interdite%
}{0525-a/4,0526-a/4,0008-a/6,0253-a/6}
\endsequence
\sequence{Division\ euclidienne}\hiddentrue
\activites{
    division_euclidienne-1%
}
\begin{exercices}
    \doublesheet{0155-c,0156-c}{0157-c,0158-c}
\end{exercices}
\endsequence
\sequence{Polygones}\hiddentrue
	\addEvalCmp{Je connais le vocabulaire associé aux quadrilatère&communiquer}
	\addEvalCmp{Je présente correctement ma(mes) feuille(s), mon écriture est lisible&transversale}
	\addEvalCmp{Je connais ma leçon&transversale}
\cours
\activites{
    polygones-1,
    quadrilateres-1%
}
\begin{exercices}
    \simplesheet{0216-c,0216-d,0217-c}
    \simplesheet{0216-e,0218-c}
    \doublesheet{0491-c}{0492-c}
\end{exercices}
\evaluationAB{
    duree=50 min,
    consignes={Sur le sujet, écrire uniquement le nom et le prénom
    et les réponses de l'exercice 1.},
    skip page id=0499%
}{0498/4,0216/6,0499/6,0219/2,0194/2}
\evaluation{
    format=DS,
    sujet=C,
    duree=50 min,
    font=sf,
    consignes={Sur le sujet, écrire uniquement le nom et le prénom.},
    skip page id=0499-a%
}{0498-a/6,0216-a/6,0499-a/6,0194-a/2}
\endsequence
\sequence{Priorit\'es\ de\ calculs}\hiddentrue
	\addEvalCmp{Je connais les priorités opératoires&calculer}
	\addEvalCmp{Je suis capable d'utiliser des astuces de calculs&calculer}
	\addEvalCmp{Je suis capable de traduire une série de calcul par une expression numérique&représenter}
	\addEvalCmp{Je redige correctement mes calculs&communiquer}
	\addEvalCmp{Je présente correctement ma(mes) feuille(s), mon écriture est lisible&transversale}
	\addEvalCmp{Je connais ma leçon&transversale}
\cours
\begin{exercices}
    \simplesheet{0334-c,0335-c}
    \simplesheet{0334-d,0335-d}
    \simplesheet{0334-e,0335-e}
\end{exercices}
\evaluation{
    format=DS,
    sujet=A,
    duree=50 min,
    calculatrice = interdite,
    consignes={Répondre directement sur la feuille.},
    skip page id=0579-c%
}{0580-c/2,0579-c/6,0579-d/6,0581-c/3,0582-c/3}
\evaluation{
    format=DS,
    sujet=B,
    duree=50 min,
    calculatrice = interdite,
    consignes={Répondre directement sur la feuille.},
    skip page id=0579-c%
}{0580-c/2,0579-c/6,0579-d/6,0581-c/3,0582-c/3}
\evaluation{
    format=DS,
    sujet=C,
    duree=50min,
    consignes={Répondre sur votre feuille.}%
}{0334-a/12,0335-a/4,0341-a/4}
\endsequence               
\sequence{Proportionnalit\'e}\hiddentrue
\cours
\activites{
    decouverte_proportionnalite-1,
    proportionnalite_linearite_additive-1,
    pourcentages-1,
    echelle-1%
}
\begin{exercices}
    \doublesheet{0390-c,0391-c}{0392-c,0393-c}
    \doublesheet{0407-d}{0408-d}
    \doublesheet{0399-c,0399-d}{0399-e}
    \doublesheet{0455-c,0456-c}{0458-c,0457-c}
    \doublesheet{0459-c,0460-c}{0461-c,0462-c}
    \doublesheet{0463-c,0464-c}{0465-c}
\end{exercices}
\questionsFlash{
    multiplication_a_trou-1,
    multiplication_a_trou-2,
    proportionnalite-1%
}
\evaluationAB{
    duree=50 min,
    consignes={Exercice 4 à faire directement sur le sujet.},
    skip page id=0407%
}{0155/4,0156/3,0158/3,0399/2,0407/3,0456/3,0463/2}
\evaluation{
    format=DS,
    sujet=C,
    duree=50 min,
    font=sf,
    consignes={Exercice 3 à faire directement sur le sujet.},
    skip page id=0407-a%
}{0155-a/4,0156-a/3,0399-a/4,0407-a/3,0456-a/3,0463-a/2}
\endsequence
\sequence{P\'erim\`etres}\hiddentrue
	\addEvalCmp{Je présente correctement ma(mes) feuille(s), mon écriture est lisible&transversale}
	\addEvalCmp{Je connais ma leçon&transversale}
	\addEvalCmp{Je suis capable de calculer le périmètre d'une figure plane&calculer}
	\addEvalCmp{Je redige correctement mes calculs&communiquer}
\cours
\activites{
    conversions_unites_longueurs-1,
    perimetre_d_un_cercle-1%
}
\begin{exercices}
    \doublesheet{0311-c,0313-c}{0314-c}
    \doublesheet{0315-c,0316-c}{0317-c}
    \doublesheet{0318-c}{0319-c,0320-c}
    \doublesheet{0099-c,0101-c}{0100-c}
    \doublesheet{0112-c}{0113-c}
    \simplesheet{0114-c}
    \doublesheet{0534-c,0535-c}{0536-c,0537-c}
\end{exercices}
\questionsFlash{%
    polygones_a_main_levee-1,
    conversion_unites_longueurs-1,
    diviseurs-1,
    conversion_unites_longueurs-2,
    multiplication_division_par_10_100_1000%
}
\evaluationAB{
    duree=50 min,
    consignes={Répondre sur votre feuille. Ne pas oublier les calculs
    et la phrase réponse.},
    endspace=15pt,
    skip page id=0311,
    calculatrice=autorisee%
}{0549/3,0548/2,0311/4,0318/3,0320/3,0100/1,0537/4,0114/2 bonus}
\evaluation{
    format=DS,
    sujet=C,
    duree=50 min,
    font=sf,
    consignes={Répondre sur votre feuille. Ne pas oublier les calculs
    et la phrase réponse.},
    endspace=15pt,
    calculatrice=autorisee%
}{0549-a/4,0548-a/6,0318-a/3,0320-a/3,0100-a/2,0537-a/2,0114-a/2 bonus}
\endsequence    
\sequence{M\'ediatrice\ d'un\ segment}\hiddentrue
\cours
\activites{
    mediatrice_equidistance-1,
    construction_mediatrice_compas-1,
    trouver_le_centre_dun_cercle-1,
    placer_centre_cercle-1,
    construction_symetrique-1%
    }
\begin{exercices}
    \doublesheet{0067-c}{0068-c}
    \simplesheet{0069-c}
    \doublesheet{0073-c}{0074-c}
    \doublesheet{0075-c}{0076-c,0077-c}
    \doublesheet{0029-c}{0030-c}
    \doublesheet{0034-c}{0035-c}
    \doublesheet{0036-c}{0037-c}
\end{exercices}
\questionsFlash{%
    calculs_priorites-2,
    calculs_priorites-3,
    calculs_priorites-4,
    calculs_priorites-5,
    axes_de_symetrie,
    conversion_unites_longueurs-4%
}
\evaluationAB{
    duree=50 min,
    consignes={Répondre sur le sujet uniquement pour l'exercice 2.},
    skip page id=0102%
}{0067/4,0069/4,0102/4,0348/4,0349/4}
\evaluation{
    duree=50 min,
    format=DS,
    sujet=C,
    consignes={Répondre sur le sujet uniquement pour l'exercice 2.},
    skip page id=0102-a%
}{0067-a/5,0069-a/5,0102-a/5,0348-a/5}
\endsequence
\sequence{Ordre\ et\ comparaison}\hiddentrue
\evaluation{
        format=DM,
        duree=À rendre le 03/01/2023,
        font=sf,
        consignes={Répondre directement sur le sujet.}
}{0222-c/6,0223-c/6,0224-c/6,0225-c/2}
\evaluationAB{
    duree=50 min,
    consignes={Les figures doivent être tracées sur le sujet.
    Les autres réponses doivent être rédigées sur la feuille.},
    skip page id=0247%
}{0222/3,0224/3,0247/4,0248/5,0249/5,0250/2 (BONUS)}
\evaluation{
    format=DS,
    sujet=C,
    duree=50 min,
    font=sf,
    consignes={Les figures doivent être tracées sur le sujet.
    Les autres réponses doivent être rédigées sur la feuille.},
    skip page id=0223-c%
}{0222-a/6,0224-a/6,0247-a/4,0248-a/4,0250-a/2 (BONUS)}
\endsequence
\sequence{Aires}\hiddentrue
	\addEvalCmp{Je présente correctement ma(mes) feuille(s), mon écriture est lisible&transversale}
	\addEvalCmp{Je connais ma leçon&transversale}
	\addEvalCmp{Je suis capable de calculer l'aire d'une figure plane&calculer}
	\addEvalCmp{Je redige correctement mes calculs&communiquer}
\cours
\activites{
    conversions_unites_aires-1,
    aire_rectangle-1,
    aire_triangle_rectangle-1%
    }
\begin{exercices}
    \doublesheet{0323-c,0324-c}{}
    \doublesheet{0543-c,0544-c}{0545-c}
    \doublesheet{0546-c}{0547-c}
    \doublesheet{0556-c,0557-c}{0558-c,0559-c}
    \doublesheet{0325-c,0326-c}{0329-c,0327-c,0328-c}
    \simplesheet{0330-c}
\end{exercices}
\questionsFlash{
    calculs_priorites_10,
    nature_des_angles-3,
    calcul_perimetre,
    aire_rectangle-1,
    aire_rectangle_triangle-1%
}
\evaluationAB{
    duree=50 min,
    consignes = {Répondre sur votre feuille.},
    calculatrice = autorisee,
    skip page id=0556%
}{0601/3,0323/4,0545/2,0556/4,0326/2,0330/4,0325/BONUS(1)}
\evaluation{
    duree = 50 min,
    format=DS,
    sujet=C,
    consignes = {Répondre sur votre feuille.},
    calculatrice = autorisee,
    skip page id=0325%
}{0323-a/4,0545-a/4,0556-a/4,0325-a/4,0330-d/4}
\endsequence
\sequence{Sym\'etrie\ axiale}\hiddentrue
	\addEvalCmp{Je présente correctement ma(mes) feuille(s), mon écriture est lisible&transversale}
	\addEvalCmp{Je sais tracer l'image d'un polygone par une symétrie axiale&représenter}
\cours
\activites{
    symetrie_axiale-1,
    symetrique_point_par_rapport_droite-1,
    construction_symetrique-1,
    construction_symetrique_figure_complexe-1,
    construction_symetrique_compas-1,
    propriete_de_conservation_symetrie_axiale-2%
}
\begin{exercices}
    \doublesheet{0041-c}{0042-c}
    \simplesheet{0028-c}
    \doublesheet{0029-c}{0030-c}
    \doublesheet{0031-c}{0032-c}
    \simplesheet{0265-c}
    \simplesheet{0266-c}
    \doublesheet{0036-c}{0035-c,0037-c}
\end{exercices}
\questionsFlash
{
    multiplier_par_05_et_025-1,
    multiplier_par_05_et_025-2,
    symetrique_oeil_nu,
    elever_au_carre-1,
    calculs_priorites-1,
    multiplication_division_par_10_100_1000%
}
\evaluation{
    format=DS,
    sujet=A,
    duree=50 min,
    consignes={},
    skip page id=0028-c%
}{0041-c/2,0028-c/3,0265-c/3,0612-c/3,0613-c/3,0614-c/3,0266-c/3}
\endsequence
\sequence{Division}\hiddentrue
\cours
\activites{%
    a_est_un_diviseur_de_b-1,
    criteres_de_divisibilite-1%
}
\begin{exercices}
    \doublesheet{0274-c,0275-c,0276-c}{0277-c,0278-c,0279-c}
    \simplesheet{0280-d}
    \doublesheet{0296-c}{0297-c}
    \doublesheet{0306-c,0307-c}{0308-c}
\end{exercices}
\evaluationAB{
    duree=50 min,
    consignes={Tous les exercices doivent être traités sur la feuille.}
}{0270/3,0271/3,0272/3,0273/3,0306/4,0296/4}
\evaluation{
    format=DS,
    sujet=C,
    duree=50 min,
    font=sf,
    consignes={Tous les exercices doivent être traités sur la feuille.}
}{0270-c/3,0271-c/3,0272-c/3,0273-c/3,0306-a/4,0296-a/4}
\endsequence             
\sequence{Prori\'et\'es\ des\ droites\ parall\`eles\ et\ }\hiddentrue
\evaluationAB{
    duree=50 min,
    consignes={Ne pas répondre sur le sujet. Les exercices peuvent-être
    traités dans l'ordre souhaité.},
    skip page id=0117%
}{0115/4,0116/4,0117/4,0118/4,0358/3,0359/3}
\evaluation{
    duree=50 min,
    format=DS,
    sujet=C,
    consignes={Répondre directement sur le sujet.}%
}{0116-a/5,0117-a/5,0115-a/5,0358-a/5}
\endsequence
\sequence{Organisation\ et\ gestion\ de\ donn\'ees}\hiddentrue
\cours
\evaluationAB{
    duree=50 min,
    consignes={Répondre directement sur le sujet.},
    multicols=\multicolstrue
}{0352-c/4,0353/4,0357/4,0354/4,0355/4,0356/4}
\evaluation{
    duree=50 min,
    format=DS,
    sujet=C,
    consignes={Répondre directement sur le sujet.},
    skip page id=0354-c%
}{0352-c/4,0353-c/4,0357-c/4,0354-c/4,0356-c/4}
\endsequence
\sequence{Solides\ et\ volumes}\hiddentrue
\cours
\endsequence
\level{5e}
\sequence{Expressions\ num\'eriques}\hiddentrue
	\addEvalCmp{Je connais ma leçon&transversale}
	\addEvalCmp{Je présente correctement ma(mes) feuille(s), mon écriture est lisible&transversale}
	\addEvalCmp{Je suis capable de traduire une série de calcul par une expression numérique&représenter}
	\addEvalCmp{Je connais les priorités opératoires&calculer}
	\addEvalCmp{Je suis capable d'utiliser des astuces de calculs&calculer}
	\addEvalCmp{Je redige correctement mes calculs&communiquer}
\cours
\activites{
    expressions_numeriques-1,
    priorites_operatoires-1,
    astuces_operatoires-1,
    expressions_numeriques-2,
    ecrire_expression_numerique_longueur-1,
    ecrire_expression_numerique_aire-1%
    }
\begin{exercices}
    \simplesheet{0360-c,0360-d}
    \simplesheet{0334-c,0334-d,0371-c}
    \simplesheet{0334-e}
    \doublesheet{0374-c,0375-c}{0372-c}
    \doublesheet{0380-c,0381-c}{0377-c}
    \doublesheet{0376-c}{0379-c}
\end{exercices}
\questionsFlash{
    operation_unique-1,
    vocabulaire_operations-1,
    operations_multiples-1,
    operations_multiples-2,
    operations_multiples-3,
    calculs_priorites-5%
    }
\evaluationAB{%
    duree=50 min,
    consignes={Ne pas répondre sur le sujet.},
    calculatrice=interdite,
    skip page id=0384%
}{0334/3,0382/6,0383/4,0384/4,0385/3}
\evaluation{
    format=DS,
    sujet=C,
    duree=50 min,
    consignes={Ne pas répondre sur le sujet.},
    calculatrice=autorisée,
    skip page id=0384-a%
}{0334-a/6,0383-a/4,0384-a/2,0385-a/4,0372-c/4}
\endsequence
\sequence{Triangles\ partie\ 1}\hiddentrue
	\addEvalCmp{Je présente correctement ma(mes) feuille(s), mon écriture est lisible&transversale}
    Je suis capable de construire un triangle en utilisant mes outils de géométrie/représenter;%
	\addEvalCmp{Je sais vérifier l'inégalité triangulaire pour savoir si un triangle est constructible&raisonner}
	\addEvalCmp{Je redige correctement mes calculs&communiquer}
\cours
\activites{
    inegalite_triangulaire_plus_court_chemin-1,
    inegalite_triangulaire_application-1,
    inconstructibilite_d_un_triangle-1,
    inegalite_triangulaire_verification-1%
}
\begin{exercices}
    \simplesheet{0454-c}
    \doublesheet{0448-c,0449-c}{0450-c, 0451-c}
    \simplesheet{0452-c}
    \doublesheet{0453-c, 0453-d}{0453-e}
\end{exercices}
\evaluationAB{
    duree=50 min,
    consignes={Ne pas répondre sur le sujet.},
    calculatrice=interdite,
    skip page id=0452,
    font=sf,
    fontsize=\large%
}{0481/4,0482/4,0452/4,0483/4,0586/2,0630/2}
\endsequence
\sequence{D\'ecouverte\ des\ nombres\ relatifs}\hiddentrue
	\addEvalCmp{Je présente correctement ma(mes) feuille(s), mon écriture est lisible&transversale}
	\addEvalCmp{Je sais comparer deux nombres relatifs&raisonner}
	\addEvalCmp{Je sais placer un point sur une droite graduée&représenter}
\cours
\activites{%
    decouverte_relatifs-1,
    comparaison_relatifs-1,
    distance_a_zero-1,
    reperage_dans_le_plan-1,
    carre_magique_relatifs-1,
    somme_relatifs-1%
}
\begin{exercices}
    \simplesheet{0436-c}
    \doublesheet{0431-c,0434-c}{0435-c}
    \doublesheet{0432-c,0433-c}{0430-c}
    \doublesheet{0437-c,0439-c}{0438-c}
    \doublesheet{0440-c}{0441-c}
    \simplesheet{0560-c,0561-c,0562-c}
    \doublesheet{0563-c}{0564-c}
    \simplesheet{0560-d,0562-d,0565-c}
\end{exercices}
\questionsFlash{%
    pourcentages-1,
    pourcentages-2,
    comparaison_relatifs-1%
}
\evaluationAB{
    duree=50 min,
    consignes={Répondre sur votre feuille.},
    calculatrice=interdite,
    skip page id=0446%
}{0442/3,0443/2,0444/4,0446/2,0447/5,0563/2,0445/2}
\evaluation{%
    format=DS,
    sujet=C,
    duree=50 min,
    font=sf,
    consignes={Répondre sur votre feuille. Répondre par des phrases.},
    calculatrice=interdite,
    skip page id=0446-a%
}{0442-a/3,0443-a/2,0444-a/6,0446-a/4,0447-a/5}
\endsequence
\sequence{Calcul\ fractionnaire}\hiddentrue
	\addEvalCmp{Je présente correctement ma(mes) feuille(s), mon écriture est lisible&transversale}
	\addEvalCmp{Je sais calculer avec les fractions&calculer}
	\addEvalCmp{Je sais utiliser des fractions pour résoudre des problèmes&raisonner}
\cours
\activites{
    simplification_de_fractions-1%
}
\begin{exercices}
    \doublesheet{0609-c}{0610-c,0611-c}
    \simplesheet{0583-c,0584-c}
    \simplesheet{0644-c,0643-c}
    \simplesheet{0645-c,0646-c}
\end{exercices}
\evaluationAB{%
    duree=50 min,
    consignes={Ne pas répondre sur le sujet.},
    calculatrice=interdite,
    skip page id=0651%
}{0647/2,0648/4,0649/3,0650/2,0651/4,0653/2,0652/3}
\evaluation{%
    format=DS,
    sujet=C,
    duree=50 min,
    consignes={Ne pas répondre sur le sujet.},
    calculatrice=autorisée%
}{0647-a/2,0648-a/4,0649-a/3,0650-a/2,0651-a/4}
\evaluation{%
    format=DS,
    sujet=D,
    duree=30 min,
    consignes={Ne pas répondre sur le sujet.},
    calculatrice=autorisée%
}{0647-c/2,0648-c/4,0649-c/3,0650-c/2,0651-c/4}
\endsequence
\sequence{Calcul\ litt\'eral}\hiddentrue
	\addEvalCmp{Je produit une expression littérale pour modéliser un problème&modéliser}
	\addEvalCmp{Je calcule avec des expressions littérales&calculer}
	\addEvalCmp{Je présente correctement mes calculs&communiquer}
\cours
\activites{
    introduction_calcul_litteral-1,
    conventions_ecritures-1,
    evaluer_une_expression_litterale-1,
    tester_une_egalite-1,
    reduire_une_expression_litterale-1,
    reduire_une_expression_litterale-2%
}
\begin{exercices}
    \doublesheet{0466-c}{0471-c}
    \doublesheet{0468-c}{0467-c}
    \doublesheet{0472-c}{0490-d}
    \doublesheet{0469-c}{0470-c}
    \doublesheet{0474-c,0475-c}{0476-c}
    \doublesheet{0477-c,0478-c}{0479-c,0480-c}
    \simplesheet{0656-c,0656-d}
    \simplesheet{0656-e,0656-f,0656-g}
    \doublesheet{0484-d,0489-d}{0473-c}
    \doublesheet{0468-d,0474-d,0478-d}{0467-d}
\end{exercices}
\evaluationAB{
    duree=50 min,
    consignes={Ne pas répondre sur le sujet.},
    calculatrice=autorisée%
    }{0474/6,0467/4,0469/4,0479/3,0471/3}
\evaluation{
        format=DS,
        sujet=C,
        font=sf,
        duree=50 min,
        consignes={Ne pas répondre sur le sujet.},
        calculatrice=autorisée%
}{0656-g/6,0474-a/6,0469-a/4,0479-a/4}
\evaluation{
        format=DS,
        sujet=D,
        font=sf,
        duree=50 min,
        consignes={Ne pas répondre sur le sujet.},
        calculatrice=autorisée%
}{0656-g/4,0474-a/2,0469-a/4,0479-a/2,0468-d/4,0467-d/4}
\evaluation{
    format=DM,
    duree=À rendre le ,
    consignes={Toutes les réponses doivent être rédigées sur une feuille
    à carreau format A4. Détailler les étapes de calcul.}%
    }{0484-c/4,0489-c/4,0490-c/2}
\endsequence
\sequence{Cas\ d'\'egalit\'e\ des\ triangles}\hiddentrue
\addEvalCmp{Je connais ma leçon&transversale}
\addEvalCmp{Je présente correctement ma(mes) feuille(s), mon écriture est lisible&transversale}
\addEvalCmp{Je démontre des égalités d'angles ou de longueurs en utilisant les cas d'égalité des triangles.&raisonner}
\cours
\activites{
    demontrer-1,
    demontrer-3,
    cas_d_egalite_CAC-1,
    cas_d_egalite_ACA-1%
}
\begin{exercices}
    \doublesheet{0494-c}{0493-c}
    \doublesheet{0673-c}{0674-c}
    \doublesheet{0501-a}{0682-c}
    \simplesheet{0502-a}
\end{exercices}
\questionsFlash{%
    criteres_de_divisibilites-1,
    fraction_en_pourcentages-1%
}
\setEvalCmp{\cmpItem0\cmpItem1\cmpItem2}
\evaluationAB{
    duree=50 min,
    consignes={Ne pas répondre sur le sujet.},
    calculatrice=autorisée,
    font=sf,
    skip page id=0501%
}{0500/2,0501/6,0502/6,0503/6}
\endsequence
\sequence{Proportionnalit\'e}\hiddentrue
	\addEvalCmp{Je connais ma leçon&transversale}
	\addEvalCmp{Je présente correctement ma(mes) feuille(s), mon écriture est lisible&transversale}
	\addEvalCmp{Je sais résoudre un problème de proportionnalité en utilisant la méthode de mon choix&raisonner}
    \addEvalCmp{Je suis capable de déterminer un nombre manquant dans un tableau de proportionnalite&calculer}
	\addEvalCmp{Je redige correctement mes calculs&communiquer}
\cours
\activites{
    a_times_x_equal_b-1,
    linearite_multiplication-1,
    proportionnalite_linearite_additive-1,
    coefficient_de_proportionnalite-1%
}
\begin{exercices}
    \doublesheet{0387-c,0389-c}{0388-c}
    \doublesheet{0390-c,0391-c}{0392-c,0393-c}
    \doublesheet{0394-c}{0395-c}
    \doublesheet{0399-c,0399-d}{0399-e,0400-c}
    \doublesheet{0401-c}{0402-c}
    \doublesheet{0407-c,0408-c,0409-c}{0410-c}
    \doublesheet{0411-c,0412-c,0413-c}{0414-c,0415-c}
    \doublesheet{0416-c,0417-c,0418-c}{0419-c,0420-c}
\end{exercices}
\questionsFlash{
    traduction_phrase_expression_numerique-1%
}
\setEvalCmp{\cmpItem2}
\evaluation{%
    format=DM,
    duree=Pour le 09/10/23,
    consignes={Répondre directement sur la feuille.},
    skip page id=0398-c%
}{0396-c/3,0397-c/3,0398-c/2,0425-c/2}
\setEvalCmp{\cmpItem0\cmpItem1\cmpItem2\cmpItem3\cmpItem4}
\evaluationAB{
    duree=50 min,
    consignes={Ne pas répondre sur le sujet.},
    calculatrice=autorisée,
    baselineskip=15pt,
    skip page id=0426%
}{0422/4,0421/3,0424/4,0417/2,0425/3,0426/4}
\evaluation{%
    format=DS,
    sujet=C,
    duree=50 min,
    font=dys,
    baselineskip=20pt,
    consignes={Ne pas répondre sur le sujet.},
    calculatrice=autorisée
}{0422-a/4,0421-a/5,0424-a/4,0417-a/4,0425-a/3}
\endsequence
\sequence{Angles\ et\ parall\'elisme}\hiddenfalse
	\addEvalCmp{Je connais le vocabulaire associé aux angles&communiquer}
	\addEvalCmp{Je démontre des égalités d'angles ou des propriétés de parallélisme&raisonner}
\cours
\activites{
    caracterisation_angulaire_du_parallelisme-1,
    caracterisation_angulaire_du_parallelisme-2%
}
\begin{exercices}
    \doublesheet{0495-c,0496-c}{0497-c}
    \doublesheet{0504-c,0505-c}{0506-c,0507-c}
    \doublesheet{0508-c}{0509-c,0510-c}
    \doublesheet{0511-c,0512-c}{0513-c}
\end{exercices}
\questionsFlash{
    vocabulaire_angles_cinquieme-1,
    reduire_expression_litterale-1,
    reduire_expression_litterale-2%
}
\evaluationAB{
    duree=50 min,
    consignes={Répondre sur votre feuille, vous pouvez annoter le sujet.},
    calculatrice=autorisée,
    skip page id=0522%
}{0521/5,0522/3,0495/2,0523/5,0524/5}
\evaluation{
    format=DS,
    sujet = C,
    font = {Comic Sans MS},
    consignes={Répondre sur votre feuille, vous pouvez annoter le sujet.},
    skip page id=0522-a%
}{0521-a/5,0522-a/3,0523-a/4,0530-c/5,0524-a/3}
\endsequence
\sequence{Triangles\ partie\ 3}\hiddenfalse
    \addEvalCmp{Je présente correctement ma(mes) feuille(s), mon écriture est lisible&transversale}
    \addEvalCmp{Je sais calculer des angles en utilisant les propriétés du cours&raisonner}
    \addEvalCmp{Je redige correctement mes calculs&communiquer}
\cours
\activites{
    somme_angles_triangle-1%
}
\begin{exercices}
    \doublesheet{0527-c,0530-c}{0528-c,0529-c,0531-c}
\end{exercices}
\questionsFlash{
    pourcentages-3,
    angle_inconnu_triangle-1%
}
\endsequence
\sequence{Division\ euclidienne}\hiddentrue
	\addEvalCmp{Je présente correctement ma(mes) feuille(s), mon écriture est lisible&transversale}
	\addEvalCmp{Je sais résoudre un problème de division&modéliser}
	\addEvalCmp{Je sais calculer des angles en utilisant les propriétés du cours&raisonner}
	\addEvalCmp{Je suis capable de déterminer si un nombre est diviseurs (multiple) d'un autre nombre&calculer}
	\addEvalCmp{Je sais utiliser le vocabulaire : multiples, diviseurs&communiquer}
\cours
\activites{
    multiples_diviseurs_schematisation-1,
    criteres_de_divisibilite-1%
}
\begin{exercices}
    \simplesheet{0538-c,0532-c}
    \doublesheet{0533-c,0539-c}{0540-c,0541-c,0542-c}
\end{exercices}
\questionsFlash{
    polygones_a_main_levee-1,
    multiplication_division_par_10_100_1000,
    pourcentages-4%
}
\evaluationAB{
    duree=50 min,
    consignes={Répondre sur votre feuille, vous pouvez annoter le sujet.},
    calculatrice=interdite%
}{0550/3,0551/3,0552/2,0553/6,0554/3,0555/3}
\evaluation{
    format=DS,
    sujet = C,
    font = sf,
    skip page id=0553-a,
    consignes={Répondre sur votre feuille excepté pour l'exercice 1, vous pouvez annoter le sujet.},
    calculatrice=autorisée%
}{0550-a/3,0551-a/3,0552-a/2,0553-a/4,0530-c/4,0529-c/2,0555-a/1}
\endsequence
\sequence{Somme\ et\ diff\'erence\ de\ nombres\ relatifs}\hiddenfalse
	\addEvalCmp{Je présente correctement ma(mes) feuille(s), mon écriture est lisible&transversale}
	\addEvalCmp{Je sais résoudre un problème de division&modéliser}
	\addEvalCmp{Je sais appliquer les règles de calcul sur les nombres relatifs&calculer}
\cours
\activites{
    % carre_magique_relatifs-1,
    % somme_relatifs-1,
    % regles_calcul_somme_relatif-1,
    soustraction_relatifs-1%
    }
\begin{exercices}
    % \simplesheet{0560-c,0561-c,0562-c}
    % \doublesheet{0563-c}{0564-c}
    \simplesheet{0560-d,0562-d,0565-c}
    \simplesheet{0566-c,0566-d,0570-c}
    \simplesheet{0571-c,0571-d}
    \simplesheet{0574-c,0575-c,0576-c,0577-c}
\end{exercices}
\questionsFlash{
    le_compte_est_bon_simple-4,
    somme_entiers_relatifs-1,
    difference_entiers_relatifs-1%
}
\evaluation{
    format=DM,
    duree=À rendre le 06/05/2024,
    consignes={Répondre directement sur la feuille.}
}{0567-c/4,0568-c/3,0569-c/4}
\evaluation{
    format=DS,
    sujet=A,
    calculatrice= interdite,
    duree = 50 min,
    consignes={Ne pas répondre sur le sujet et détailler les étapes
    de calcul.},
    skip page id=0576-d%
}{0563-d/4,0573-d/3,0574-d/3,0575-d/3,0576-d/3,0577-d/3,0578-c/3}
\endsequence
\sequence{Symetrie\ centrale}\hiddentrue
	\addEvalCmp{Je présente correctement ma(mes) feuille(s), mon écriture est lisible&transversale}
	\addEvalCmp{Je sais résoudre un problème de division&modéliser}
	\addEvalCmp{Je sais appliquer les règles de calcul sur les nombres relatifs&calculer}
\cours
\questionsFlash{
    somme_et_difference_entiers_relatifs-1%
}
\begin{exercices}
    \doublesheet{0595-c}{0596-c}
    \doublesheet{0597-c,0600-c}{0598-c}
    \simplesheet{0599-c}
\end{exercices}
\evaluation{%
    format=DS,
    sujet=C,
    duree=50 min,
    consignes={Répondre directement sur le sujet. Les construction doivent être effectuées au crayon à papier.},
    font=sf,
    skip page id=0598-c%
}{0596-c/3,0598-c/3,0607-c/4,0603-c/4,0605-c/4,0606-c/4,0604-c/4}
\endsequence
\level{3e}
\sequence{Arithm\'etique}\hiddentrue
	\addEvalCmp{Je présente correctement ma(mes) feuille(s), mon écriture est lisible&transversale}
	\addEvalCmp{Résoudre un problème d'arithmétiques&modéliser}
	\addEvalCmp{Je redige correctement mes calculs&communiquer}
\cours
\activites{
    juniper_green-1,
    juniper_green-2,
    crible_d_eratostene-1%
}
\questionsFlash{
    diviseurs_et_multiples-1,
    diviseurs_et_multiples-2,
    notations_geometrie-4,
    diviseurs-2,
    arithmetique-1%
}
\begin{exercices}
    \doublesheet{0615-c,0616-c}{0617-c}
    \doublesheet{0618-c,0619-c}{0620-c}
    \simplesheet{0137-a,0138-a}
    \doublesheet{0622-c,0623-c}{0624-c}
    \simplesheet{0160-c}
\end{exercices}
\evaluationAB{%
    duree=50 min,
    calculatrice=autorisée,
    skip page id=0621,
    presentation=2%
}{0162/5,0161/5,0621/5,0159/3}
\evaluation{%
    format=DS,
    sujet=C,
    duree=50 min,
    font=sf,
    calculatrice=autorisée,
    presentation=2%
}{0162-b/6,0621-c/6,0161-c/6}
\endsequence
\sequence{Triangles\ semblables}\hiddentrue
	\addEvalCmp{Je présente correctement ma(mes) feuille(s), mon écriture est lisible&transversale}
	\addEvalCmp{Résoudre un problème de géométrie&raisonner}
    \addEvalCmp{Je rédige correctement mes démonstrations&communiquer}
\cours
\begin{exercices}
\simplesheet{0144-c}
\simplesheet{0637-c,0626-c}
\doublesheet{0627-c}{0628-c}
\doublesheet{0181-c}{0182-c}
\end{exercices}
\evaluationAB{%
    duree=50 min,
    calculatrice = autorisée,
    presentation=2%
}{0636/4,0626/4,0182/4,0635/6}
\evaluation{%
    format=DS,
    sujet=C,
    font=sf,
    duree=50 min,
    calculatrice = autorisée,
    skip page id=0182-c,
    presentation=2%
}{0636-c/4,0182-c/6,0635-c/8}
\endsequence
\sequence{Calcul\ litt\'eral}\hiddentrue
	\addEvalCmp{Je sais calculer des expressions littérales&calculer}
\cours
\questionsFlash{%
    distributivite-1%
}
\activites{
    simple_distributivite-1,
    double_distributivite-1,
    identites_remarquables-1%
    }
\begin{exercices}
    \doublesheet{0631-c,0632-c}{0633-c}
    \doublesheet{0634-c}{0185-c}
    \simplesheet{0191-c,0192-c}
\end{exercices}
\evaluationAB{
    duree=50 min,
    consignes={Ne pas répondre sur le sujet.},
    calculatrice=autorisée,
    presentation=2,
    skip page id=0641%
}{0638/4,0639/4,0050/4,0640/2,0641/2,0642/2 Bonus}
\evaluation{
    format=DS,
    sujet=C,
    font=sf,
    duree=50 min,
    consignes={Ne pas répondre sur le sujet.},
    calculatrice=autorisée,
    presentation=2%
}{0191-a/8,0183-a/6,0201-a/4}
\endsequence
\sequence{Statistiques}\hiddentrue
\cours
\begin{exercices}
    \afoursheet{0212-c}
    \simplesheet{0213-c}
    \afoursheet{0214-c}
    \afoursheet{0215-c}
\end{exercices}
\endsequence
\sequence{Equations}\hiddentrue
\cours
\diaporamas{%
    introduction_aux_equations%
}
\begin{exercices}
    \doublesheet{0654-a,0654-b}{0654-c,0654-d}
    \doublesheet{0654-e,0654-f}{0654-g,0654-h}
    \doublesheet{0655-a,0655-b}{0655-c,0655-d}
    \doublesheet{0655-g, 0655-h}{0657-c}
    \doublesheet{0655-i, 0655-j}{0658-c}
    \simplesheet{0655-e,0655-f}
    \simplesheet{0229-c,0230-c,0231-c}
    \simplesheet{0659-c}
\end{exercices}
\evaluationAB{
    duree=50 min,
    consignes={Ne pas répondre sur le sujet.},
    skip page id=0664,
    calculatrice=autorisée,
    presentation=2%
}{0662/5,0664/5,0661/5,0654-b/3}
\evaluation{
    format=DS,
    sujet=C,
    duree=50 min,
    consignes={Ne pas répondre sur le sujet.},
    skip page id=0659-c,
    calculatrice=autorisée,
    presentation=2%
}{0662-c/5,0659-c/5,0660-c/5,0654-c/3}
\endsequence
\sequence{Fonctions}\hiddentrue
\addEvalCmp{Je connais les différents modes de représentation d'une fonction&représenter}
\cours
\activites{
    fonction_decouverte_machines-1,
    fonction_representation_tableau-2,
    fonction_representation_algebrique_tableau-1%
}
\begin{exercices}
    \doublesheet{0667-c}{0668-c}
    \doublesheet{0669-c,0669-d}{0670-c,0671-c,0672-c}
    \simplesheet{0237-c,0238-c}
    \simplesheet{0239-c,0240-c}
    \afoursheet{0675-c}
    \simplesheet{0681-c}
    \simplesheet{0261-c}
    \doublesheet{0251-c}{0252-c}
    \afoursheet{0683-c}
    \simplesheet{0685-c}
    \afoursheet{0238-c,0240-c,0261-c}
\end{exercices}
\questionsFlash{%
    dnb_programme_de_calcul-1,
    criteres_de_divisibilites-1,
    fraction_en_pourcentages-1,
    calculer_image_fonction-1%
}
\setEvalCmp{\cmpItem0}
\evaluationAB{
    duree=50 min,
    consignes={Ne pas répondre sur le sujet, excepté pour
    la question {\bfseries 5} de l'exercice 3.},
    calculatrice=autorisée,
    presentation=2,
    skip page id=0641%
}{0238/4,0240/4,0261/4,0684/4}
\evaluationAB{
    format=CC,
    duree=30min,
    consignes={Écrire directement sur le sujet.},
    skip page id=0260%
}{0257/4,0258/4,0259/4,0260/4,0261/4}
\evaluation{
    format=DM,
    duree={À rendre le 06/03/2023},
    consignes={}
}{0282-c/10,0283-c/10}
\endsequence
\sequence{Trigonom\'etrie}\hiddentrue
\addEvalCmp{Je connais mes formules trigonométriques.&transversale}
\addEvalCmp{Je calcule des mesures d'angles en utilisant les fonctions trigonométriques&calculer}
\addEvalCmp{Je rédige correctement mes démonstrations géométriques.&communiquer}
\cours
\activites{
    fonction_cos_et_sin-1,%
    decouverte_trigonometrie-1,%
    trigonometrie_calculer_un_angle-1%
}
\begin{exercices}
    \simplesheet{0687-c}
    \doublesheet{0293-c,0294-c,0292-c}{0295-c}
    \simplesheet{0688-c}
    \afoursheet{0689-c}
    \afoursheet{0690-c,0691-c}
    \afoursheet{0692-c}
\end{exercices}
\setEvalCmp{\cmpItem0\cmpItem1}
\evaluationAB{
    format=CC,
    sujet=A,
    duree=20 min,
    consignes={Répondre directement sur le sujet.},
    skip page id=0300%
}{0298/2,0299/2,0300/3,0301/3}
\setEvalCmp{\cmpItem0\cmpItem1\cmpItem2}
\evaluationAB{
    format=DS,
    sujet=A,
    duree=50 min,
    presentation=2,
    skip page id=0302,
    answerspaces={{id=0300,type=rectangle}}%
}{0300/3,0299/3,0304/3,0305/3,0302/6}
\endsequence
\sequence{R\'eciproque\ du\ th\'eor\`eme\ de\ Thal\`es}\hiddentrue
    \addEvalCmp{Je présente correctement ma(mes) feuille(s), mon écriture est lisible&transversale}
    \addEvalCmp{Je sais résoudre un problème de géométrie&raisonner}
    \addEvalCmp{Je rédige correctement mes démonstrations&communiquer}
    \cours
    \activites{
        reciproque_thales-1%
        }
    \begin{exercices}
        \doublesheet{0309-c}{0310-c}
        \simplesheet{0709-c}
        \afoursheet{0711-c}
        \afoursheet{0715-c}
    \end{exercices}
    \evaluationAB{
        duree=50 min,
        presentation=2,
        consignes={Ne pas répondre sur le sujet.},
        skip page id=0310%
    }{0712/3,0309/5,0310/5,0710/5}
    \evaluation{
        format=DM,
        duree=À rendre le 10/05/2023,
        consignes={Répondre sur une feuille A4.},
        skip page id=0336-c%
    }{0336-c/5,0337-c/5}
\endsequence
\sequence{Transformations}\hiddentrue
    \cours
    \activites{
        decouverte_homothetie-1%
    }
    \endsequence
\sequence{Fonctions\ lin\'eaires\ et\ affines}\hiddentrue
\cours
\activites{
    fonctions_lineaires-1,
    fonctions_affines-1%
}
\begin{exercices}
    \simplesheet{0338-c,0339-c,0340-c}
    \simplesheet{0338-d,0339-d,0340-d}
\end{exercices}
\evaluationAB{
    duree=50 min,
    consignes={La présentation et la qualité de la rédaction
    comptera pour 2 points dans la note finale.},
    skip page id=0346%
}{0339/4,0345/4,0346/4,0347/6}
\evaluation{
    duree=50 min,
    format=DS,
    sujet=D,
    consignes={La présentation et la qualité de la rédaction
    comptera pour 2 points dans la note finale.}%
}{0339-d/4,0345-d/4,0346-a/4,0347-a/6}
\endsequence
\end{document}